\paragraph{Realizability} The second part of fundamental theorem provide guarantee for realizability, i.e., a trace can be produce by execution of well-typed term. We say that a trace \emph{realizes} a buffer $\{\overline{m_i}\}$ when
      it contains all messages in this buffer, i.e.,
      {\small$\alpha_1 \listconcat [m_1] \listconcat ... [m_{n}] \listconcat
      \alpha_{n+1}$}. We also generalize this idea to automata.

\newcommand{\realz}{\lesssim}

\begin{definition}[Trace realize buffer]\label{def:tr-buffer} A trace $\alpha$ realizes buffer $\{\overline{m_i}\}$ when it contains all messages in this buffer, i.e., 
{\small$\alpha = \alpha_1 \listconcat [m_1] \listconcat ... [m_{n}] \listconcat \alpha_{n+1}$}, denoted as $\beta \realz \alpha$. 
\end{definition}

\begin{definition}[Automata realize buffer]\label{def:a-buffer} A automata $F$ realizes the buffer $\beta$ iff $\exists \alpha \in \denot{F}. \beta \realz \alpha$, denoted as $\beta \realz F$.
\end{definition}

We now prove a stronger theorem than the second part of the fundamental theorem, where we additionally require that the message buffer after the execution of a well-typed term can be realized by the prophecy automata of the \PAT{}:

\begin{theorem}[Realizability]\label{theorem:realizable} Given a well-formed handler specification $\Delta$, A well typed program $e$ at least realize one trace:
 {\small\begin{align*} 
&\Gamma;\Delta\;\Theta \vdash e : \rg{H}{A}{F} \impl 
\\&\qquad \forall \sigma \in
    \denotation{\Gamma}. \forall \alpha_h \in \denot{\sigma(H)}.\forall \beta \in \denot{\Theta}. \exists \alpha \in \denot{\sigma(A)}.\exists \beta'. \msteptr{\alpha_h}{(\beta, e)}{\alpha}{(\beta', ())} \land \beta' \realz \sigma(F)
\end{align*}}
\end{theorem}
\begin{proof} We proceed by induction over our type judgment
  $\Gamma;\Delta;\Theta \vdash e : \tau$, which consists of the
  following $8$ cases:
\begin{enumerate}[label=Case]

\item:
{\footnotesize
\mprooftr{87mm}
{
$\Delta(\eff{op}) =\langle \mykw{gen}\;\tau, \Theta' \rangle$ \quad
$\Gamma \vdash \tau <: \overline{x_i{:}t_i}\sarr
\rg{H}{\lastA\msg{op}{\phi}}{A \seqA F}$ \quad
$\forall i. \Gamma \vdash v_i : t_i$ \quad
 $\Gamma; \Delta; \Theta \cup \Theta' \vdash e : \rg{H \seqA \lastA\msg{op}{\phi\msubst{x_i}{v_i}}}{A}{F}$
 }
{TGen}
{$\Gamma; \Delta; \Theta \vdash \zgen{op}{\overline{v_i}}{e} : \rg{H}{\lastA\msg{op}{\phi\msubst{x_i}{v_i}} \seqA A}{F}$}
\ \\ \ \\}
This rule assume that $e \equiv \zgen{op}{\overline{v}}{e}, \tau \equiv \rg{H}{\lastA\msg{op}{\phi\msubst{x}{v}} \seqA A}{F}$, thus we need to prove
{\footnotesize\begin{align*}
\forall \sigma \in
    &\denotation{\Gamma}. \forall \alpha_h \in \denot{\sigma(H)}.\forall \beta \in \denot{\Theta}. 
    \\&\qquad \exists \alpha \in \denot{\sigma(\lastA\msg{op}{\phi\msubst{x_i}{v_i}} \seqA A)}.\exists \beta'. \msteptr{\alpha_h}{(\beta, \zgen{op}{\overline{v}}{e})}{\alpha}{(\beta', ())} \land \beta' \realz \sigma(F)
\end{align*}}\noindent
From the induction hypothesis and the precondition of this rule, we have
{\footnotesize\begin{alignat}{2}
\setcounter{equation}{0}
&\Delta(\eff{op}) =\langle \mykw{gen}\;\tau, \Theta' \rangle \ptag{assumption} \label{real:1-1} \\
&\Gamma \vdash \tau <: \overline{x{:}t}\sarr \rg{H}{\lastA\msg{op}{\phi}}{A \seqA F} \ptag{assumption} \label{real:1-2} \\
&\forall i. \Gamma \vdash v_i : t_i \ptag{assumption} \label{real:1-3} \\
&\Gamma ~|~ \Theta \cup \Theta' \vdash e : \rg{H \seqA \lastA\msg{op}{\phi\msubst{x}{v}}}{A}{F} \ptag{assumption} \label{real:1-4} \\
&\forall \sigma \in \denot{\Gamma}.\forall \alpha_h \in \denot{\sigma(H \seqA \lastA\msg{op}{\phi\msubst{x}{v}})}. \notag
\\&\quad \forall \beta \in \denot{\Theta \cup \Theta'}.\exists \alpha \in\denot{\sigma(A)}. \exists \beta'. \msteptr{\alpha_h}{(\beta, e)}{\alpha}{(\beta', ())} \land \beta' \realz \sigma(F)
\ptag{induction hypothesis} \label{real:1-5} \\
&\forall i. \forall \sigma \in \denot{\Gamma}. \sigma(v_i) \in  \sigma(\denot{t_i}) \ptag{\ref{real:1-3} and Lemma~\ref{theorem:pure-fundamental-ap}} \label{real:1-6} \\
&\forall \sigma \in \denot{\Gamma}. \sigma(\phi)\msubst{x_i}{v_i} \ptag{Lemma~\ref{lemma:built-in-typing-sub}, \ref{real:1-1}, \ref{real:1-2}, and \ref{real:1-3}} \label{real:1-7}
\end{alignat}}\noindent
From now, we consider each $\sigma \in \denot{\Gamma}$, $\beta \in \denot{\Theta}$, and $\alpha_h \in \denot{H}$ and try to prove the subgoal of this case: 
{\footnotesize\begin{align*}
    \exists \alpha \in\denot{\sigma(\lastA\msg{op}{\phi\msubst{x_i}{v_i}} \seqA A)}. \exists \beta'. \msteptr{\alpha_h}{(\beta, \zgen{op}{\overline{v}}{e})}{\alpha}{(\emptyset, ())} \land \beta' \realz \sigma(F)
\end{align*}}\noindent
Then we have
{\footnotesize\begin{alignat}{2}
&\sigma \in \denot{\Gamma} \land \beta \in \denot{\Theta} \land \alpha_h \in \denot{\sigma(H)} \ptag{assumption} \label{real:1-8}
\\& [\zevent{op}{\overline{\sigma(v_i)}}] \in \denot{\lastA\msg{op}{\sigma(\phi)\msubst{x_i}{v_i}}} \ptag{lemma~\ref{lemma:lastA-denot} and \ref{real:1-7}} \label{real:1-9}
\\& \alpha_h \listconcat [\zevent{op}{\overline{\sigma(v_i)}}] \in \denot{\sigma(H \seqA \lastA\msg{op}{\phi\msubst{x_i}{v_i}})} \ptag{lemma~\ref{lemma:seqA-denot}, \ref{real:1-8}, and \ref{real:1-9}} \label{real:1-10}
\end{alignat}}\noindent
According to the well-formed type context (Lemma~\ref{lemma:built-in-typing-sub}), \ref{real:1-1}, \ref{real:1-2}, \ref{real:1-8}, we have
{\footnotesize\begin{alignat}{2}
& \exists \beta_{\eff{op}}. \beta_{\eff{op}} \land \denot{\Theta'} \land \beta_{\eff{op}} \realz \sigma(A\seqA F) \land \alpha_h \vDash \zevent{op}{\overline{c}} \Downarrow \beta_{\eff{op}} \ptag{Lemma~\ref{lemma:built-in-typing-sub}} \label{real:1-11}
\\& \beta \cup \beta_{\eff{op}} \in  \denot{ \Theta \cup \Theta'} \ptag{Lemma~\ref{lemma:buf-partition} and \ref{real:1-11}} \label{real:1-12}
\end{alignat}}\noindent
Now, we can apply hypothesis~\ref{real:1-5} with
{\footnotesize\begin{align*}
    &\sigma \mapsto \sigma \quad \alpha_h \mapsto \alpha_h \listconcat [\zevent{op}{\overline{\sigma(v_i)}}] \quad \beta \mapsto \beta \cup \beta_{\eff{op}}
\end{align*}}\noindent
Then we have
{\footnotesize\begin{alignat}{2}
&\exists \alpha \in\denot{\sigma(A)}. \exists \beta'. \msteptr{\alpha_h}{(\beta \cup \beta_{\eff{op}}, e)}{\alpha}{(\beta', ())} \land \beta' \realz \sigma(F) \ptag{hypothesis~\ref{real:1-5} with \ref{real:1-8}, \ref{real:1-10}, and \ref{real:1-11}} \label{real:1-13}
\\& \alpha \listconcat [\zevent{op}{\overline{\sigma(v_i)}}] \in \denot{\sigma(\lastA\msg{op}{\phi\msubst{x_i}{v_i}} \seqA A)} \ptag{lemma~\ref{lemma:seqA-denot}, \ref{real:1-9}, and \ref{real:1-13}} \label{real:1-14}
\end{alignat}}\noindent
With help of hypothesis \ref{real:1-13} and \ref{real:1-14}, we can instantiate the existential quantified variables as $\alpha \mapsto [\zevent{op}{\overline{\sigma(v_i)}}] \listconcat \alpha, \beta' \mapsto \beta'$, and we need to prove
{\footnotesize\begin{align*}
    \msteptr{\alpha_h}{(\beta, \zgen{op}{\overline{v}}{e})}{[\zevent{op}{\overline{\sigma(v_i)}}] \listconcat \alpha}{(\beta', ())}
\end{align*}}\noindent
where
{\footnotesize\begin{alignat}{2}
& \alpha_h \vDash \zevent{op}{\overline{c}} \Downarrow \beta_{\eff{op}} \ptag{hypothesis \ref{real:1-11}} \label{real:1-15}
\\& \msteptr{\alpha_h}{(\beta \cup \beta_{\eff{op}}, e)}{\alpha}{(\beta', ())} \ptag{hypothesis \ref{real:1-13}} \label{real:1-16}
\\& \msteptr{\alpha_h}{(\beta, \zgen{op}{\overline{v}}{e})}{ [\zevent{op}{\overline{\sigma(v_i)}}] \listconcat \alpha}{(\beta', ())} \ptag{\textsc{StGen}, \ref{real:1-15}, and \ref{real:1-16}} \label{real:1-17}
\end{alignat}}\noindent
which is sufficient to prove the subgoal in this case.

\item:
{\footnotesize
\mprooftr{80mm}
{
  $\Delta(\eff{op}) =\langle \mykw{obs}\;\tau, \Theta'\rangle$ \quad
$\Gamma \vdash \tau <: \overline{x_i{:}t_i}\sarr
\rg{H}{\lastA\msgB{op}{\overline{y}}{\phi \land \overline{y = x}}}{A \seqA F}$ \quad
 $\Gamma, \overline{x{:}t}; \Delta; \Theta \cup \Theta' \vdash e : \rg{H \seqA \lastA\msgB{op}{\overline{y}}{\phi \land \overline{y = x}}}{A}{F}$
 }
{TObs}
{$\Gamma; \Delta; \{\eff{op}\} \cup \Theta \vdash \zobs{\overline{x}}{op}{e} : \rg{H}{\lastA\msg{op}{\phi}\seqA A}{F}$}
\ \\ \ \\}
This rule assume that $e \equiv \zobs{\overline{x}}{op}{e}, \tau \equiv\rg{H}{\lastA\msg{op}{\phi} \seqA A}{F}, \Theta \equiv \{\eff{op}\} \cup \Theta$, thus we need to prove
{\footnotesize\begin{align*}
\forall \sigma \in
    &\denotation{\Gamma}. \forall \alpha_h \in \denot{\sigma(H)}.\forall \beta \in \denot{\{\eff{op}\} \cup \Theta}. 
    \\&\qquad \exists \alpha \in \denot{\sigma(\lastA\msg{op}{\phi} \seqA A)}.\exists \beta'. \msteptr{\alpha_h}{(\beta, \zobs{\overline{x}}{op}{e})}{\alpha}{(\beta', ())} \land \beta' \realz \sigma(F)
\end{align*}}\noindent
From the induction hypothesis and the precondition of this rule, we have
{\footnotesize\begin{alignat}{2}
\setcounter{equation}{0}
&\Delta(\eff{op}) =\langle \mykw{obs}\;\tau, \Theta' \rangle \ptag{assumption} \label{real:2-1} \\
&\Gamma \vdash \tau <: \overline{x_i{:}t_i}\sarr\rg{H}{\lastA\msgB{op}{\overline{y}}{\phi \land \overline{y = x}}}{A \seqA F} \ptag{assumption} \label{real:2-2} \\
&\Gamma, \overline{x{:}t}; \Delta; \Theta \cup \Theta' \vdash e : \rg{H \seqA \lastA\msgB{op}{\overline{y}}{\phi \land \overline{y = x}}}{A}{F} \ptag{assumption} \label{real:2-3} \\
&\forall \sigma \in \denot{\Gamma, \overline{x{:}t}}.\forall \alpha_h \in \denot{\sigma(H \seqA \lastA\msgB{op}{\overline{y}}{\phi \land \overline{y = x}})}. \notag
\\&\quad \forall \beta \in \denot{\Theta \cup \Theta'}.\exists \alpha \in\denot{\sigma(A)}. \exists \beta'. \msteptr{\alpha_h}{(\beta, e)}{\alpha}{(\beta', ())} \land \beta' \realz \sigma(F)
\ptag{induction hypothesis} \label{real:2-4}
\end{alignat}}\noindent
From now, we consider each $\sigma\msubst{x_i}{v_i} \in \denot{\Gamma, \overline{x{:}t}}$, $\beta \cup \zevent{op}{\overline{\sigma(v_i)}} \in \denot{\{\eff{op}\} \cup \Theta}$, and $\alpha_h \in \denot{H}$ and try to prove the subgoal of this case: 
{\footnotesize\begin{align*}
    \qquad \exists \alpha \in \denot{\sigma(\lastA\msg{op}{\phi} \seqA A)}.\exists \beta'. \msteptr{\alpha_h}{(\beta, \zobs{\overline{x}}{op}{e})}{\alpha}{(\beta', ())} \land \beta' \realz \sigma(F)
\end{align*}}\noindent
Then we have
{\footnotesize\begin{alignat}{2}
&\sigma \in \denot{\Gamma} \land \alpha_h \in \denot{\sigma(H)} \land \beta \in \denot{\Theta} \land \forall i. \sigma(v_i) \in\denot{\sigma(t_i)} \ptag{assumption} \label{real:2-5}
\\& [\zevent{op}{\overline{\sigma(x)}}] \in \denot{\sigma(\lastA\msgB{op}{\overline{y}}{\phi \land \overline{y = x}})} \ptag{lemma~\ref{lemma:lastA-denot}} \label{real:2-6}
\\& \alpha_h \listconcat [\zevent{op}{\overline{\sigma(x)}}] \in \denot{\sigma(H \seqA \lastA\msgB{op}{\overline{y}}{\phi \land \overline{y = x}})} \ptag{lemma~\ref{lemma:seqA-denot}, \ref{real:2-5}, and \ref{real:2-6}} \label{real:2-7}
\end{alignat}}\noindent
According to the well-formed type context (Lemma~\ref{lemma:built-in-typing-sub}), \ref{real:2-1}, \ref{real:2-2}, and \ref{real:2-7}, we have
{\footnotesize\begin{alignat}{2}
& \exists \beta_{\eff{op}}. \beta_{\eff{op}} \land \denot{\Theta'} \land \beta_{\eff{op}} \realz \sigma(A\seqA F) \land \alpha_h \vDash \zevent{op}{\overline{\sigma(v_i)}} \Downarrow \beta_{\eff{op}} \ptag{Lemma~\ref{lemma:built-in-typing-sub}} \label{real:2-8}
\\& \beta \cup \beta_{\eff{op}} \in  \denot{ \Theta \cup \Theta'} \ptag{Lemma~\ref{lemma:buf-partition} and \ref{real:2-8}} \label{real:2-9}
\end{alignat}}\noindent
Now, we can apply hypothesis~\ref{real:2-4} with
{\footnotesize\begin{align*}
    &\sigma \mapsto \sigma\msubst{x_i}{v_i} \quad \alpha_h \mapsto \alpha_h \listconcat [\zevent{op}{\overline{\sigma(v_i)}}] \quad \beta \mapsto \beta \cup \beta_{\eff{op}}
\end{align*}}\noindent
Then we have
{\footnotesize\begin{alignat}{2}
&\exists \alpha \in\denot{\sigma(A)}. \exists \beta'. \msteptr{\alpha_h}{(\beta \cup \beta_{\eff{op}}, e)}{\alpha}{(\beta', ())} \land \beta' \realz \sigma(F) \ptag{hypothesis~\ref{real:2-4} with \ref{real:2-5}, \ref{real:2-7}, and \ref{real:2-9}} \label{real:2-10}
\\& \alpha \listconcat [\zevent{op}{\overline{\sigma(v_i)}}] \in \denot{\sigma(\lastA\msgB{op}{\overline{y}}{\phi \land \overline{y = x}} \seqA A)} \ptag{lemma~\ref{lemma:seqA-denot}, \ref{real:1-6}, and \ref{real:1-10}} \label{real:1-11}
\end{alignat}}\noindent
With help of hypothesis \ref{real:1-10} and \ref{real:1-11}, we can instantiate the existential quantified variables as $\beta \mapsto \{ \zevent{op}{\overline{\sigma(v_i)}} \}\cup \beta, \alpha \mapsto [\zevent{op}{\overline{\sigma(v_i)}}] \listconcat \alpha$, and we need to prove
{\footnotesize\begin{align*}
    \msteptr{\alpha_h}{(\beta, \zgen{op}{\overline{v}}{e})}{[\zevent{op}{\overline{\sigma(v_i)}}] \listconcat \alpha}{(\beta', ())}
\end{align*}}\noindent
where
{\footnotesize\begin{alignat}{2}
& \alpha_h \vDash \zevent{op}{\overline{c}} \Downarrow \beta_{\eff{op}} \ptag{hypothesis \ref{real:2-8}} \label{real:1-12}
\\& \msteptr{\alpha_h}{(\beta \cup \beta_{\eff{op}}, e)}{\alpha}{(\beta', ())} \ptag{hypothesis \ref{real:1-10}} \label{real:1-13}
\\& \msteptr{\alpha_h}{(\beta, \zgen{op}{\overline{v}}{e})}{ [\zevent{op}{\overline{\sigma(v_i)}}] \listconcat \alpha}{(\beta', ())} \ptag{\textsc{StGen}, \ref{real:1-12}, and \ref{real:1-13}} \label{real:1-14}
\end{alignat}}\noindent
which is sufficient to prove the subgoal in this case.

\item:
{\footnotesize
\mprooftr{35mm}
{}
{TRet}
{$\Gamma; \Delta; \emptyset \vdash () : \rg{H}{\epsTL}{F}$}
\ \\ \ \\}
This rule assume that $\Theta \equiv \emptyset, e \equiv (), \tau \equiv \rg{H}{\epsTL}{F}$, thus we need to prove
{\footnotesize\begin{align*}
\forall \sigma \in
    \denotation{\Gamma}. \forall \alpha_h \in \denot{\sigma(H)}.\forall \beta \in \denot{\Theta}. \exists \alpha \in \denot{\sigma(\epsTL)}.\exists \beta'. \msteptr{\alpha_h}{(\beta, ())}{\alpha}{(\beta', ())} \land \beta' \realz \sigma(F)
\end{align*}}\noindent
Note that the denotation of empty capability only contains an empty buffer, also only empty trace $[]$ is in the denotation of $\sigma(\epsTL)$. Thus, we can instantiate $\beta'$ as $\emptyset$ and prove $\msteptr{\alpha_h}{(\emptyset, ())}{[]}{(\emptyset, ())}$, which immediate holds.

\item:
{\footnotesize
\mprooftr{50mm}
{
$\Gamma; \Delta; \Theta \vdash e_1 : \rg{H}{A}{F}$ \quad $\Gamma; \Delta; \Theta \vdash e_2 : \rg{H}{A}{F}$
 }
{TChoice}
{$\Gamma; \Delta; \Theta \vdash e_1 \oplus e_2 : \rg{H}{A}{F}$}
\ \\ \ \\}
This rule assumes that $e \equiv e_1 \oplus e_2$, thus we need to prove
{\footnotesize\begin{align*}
\forall \sigma \in
    \denotation{\Gamma}. \forall \alpha_h \in \denot{\sigma(H)}.\forall \beta \in \denot{\Theta}. \exists \alpha \in \denot{\sigma(A)}.\exists \beta'. \msteptr{\alpha_h}{(\beta, e_1 \oplus e_2)}{\alpha}{(\beta', ())} \land \beta' \realz \sigma(F)
\end{align*}}\noindent
From the inductive hypothesis of this case, we know
{\footnotesize\begin{align*}
    & \forall \sigma \in
    \denotation{\Gamma}. \forall \alpha_h \in \denot{\sigma(H)}.\forall \beta \in \denot{\Theta}. \exists \alpha \in \denot{\sigma(A)}.\exists \beta'. \msteptr{\alpha_h}{(\beta, e_1)}{\alpha}{(\beta', ())} \land \beta' \realz \sigma(F)
\end{align*}}\noindent
We also know $\msteptr{\alpha_h}{(\beta, e_1 \oplus e_2)}{[]}{(\beta, e_1)}$ from \textsc{StChoice}, Then it is sufficient to prove the subgoal of this case.

\item:
{\footnotesize
\mprooftr{60mm}
{$\Gamma, z\hasoty{unit}{\phi};\Delta;\Theta \vdash e : \rg{H}{A}{F}$ \quad $z$ is fresh}
{TAssume}
{$\Gamma;\Delta;\Theta  \vdash \zassume{\phi}{e} : \rg{H}{A}{F}  $}
\ \\ \ \\}
This rule assume that $e \equiv \zassume{\phi}{e}$, thus we need to prove
{\footnotesize\begin{align*}
\forall \sigma \in
    \denotation{\Gamma}. \forall \alpha_h \in \denot{\sigma(H)}.\forall \beta \in \denot{\Theta}. \exists \alpha \in \denot{\sigma(A)}.\exists \beta'. \msteptr{\alpha_h}{(\beta, \zassume{\phi}{e})}{\alpha}{(\beta', ())} \land \beta' \realz \sigma(F)
\end{align*}}\noindent
From the inductive hypothesis of this case, we know
{\footnotesize\begin{align*}
    & \forall \sigma \in \denot{\Gamma, z\hasoty{unit}{\phi}}.\forall \alpha_h \in \denot{\sigma(H)}.\forall \beta \in \denot{\Theta}. \exists \alpha \in \denot{\sigma(A)}.\exists \beta'. \msteptr{\alpha_h}{(\beta, e)}{\alpha}{(\beta', ())} \land \beta' \realz \sigma(F)
\end{align*}}\noindent
Since $z$ is a fresh variable, then we have
{\footnotesize\begin{align*}
    &\forall \sigma, \sigma \in \denot{\Gamma, z\hasoty{unit}{\phi}} \impl \exists \sigma'. \sigma'[\overline{z \mapsto ()}] = \sigma. \sigma'(e) \in \denot{\sigma'(\tau)}
\end{align*}}\noindent
Moreover, according to the definition of type context denotation,
{\footnotesize\begin{align*}
    &\forall \sigma, \sigma \in \denot{\Gamma} \land \sigma(\phi) \iff \sigma[\overline{z \mapsto ()}] \in \denot{\Gamma, z\hasoty{unit}{\phi}}
\end{align*}}\noindent
Now, we just need to show $\zassume{\phi}{e}$ can reduced into $e$ without add new effect, which is can be proved by \textsc{StAssume} and $\sigma(\phi)$. Then the proof immediate holds in this case.

\item:
{\footnotesize
\mprooftr{45mm}
{$\Gamma;\Delta;\Theta \vdash e : \rg{H}{A}{F}$ \quad
$\Gamma \vdash () : \nuot{unit}{\phi}$}
{TAssert}
{$\Gamma;\Delta;\Theta \vdash \zassert{\phi}{e} : \rg{H}{A}{F}  $}
\ \\ \ \\}
This rule assume that $e \equiv \zassert{\phi}{e}$, thus we need to prove
{\footnotesize\begin{align*}
\forall \sigma \in
    \denotation{\Gamma}. \forall \alpha_h \in \denot{\sigma(H)}.\forall \beta \in \denot{\Theta}. \exists \alpha \in \denot{\sigma(A)}.\exists \beta'. \msteptr{\alpha_h}{(\beta, \zassert{\phi}{e})}{\alpha}{(\beta', ())} \land \beta' \realz \sigma(F)
\end{align*}}\noindent
From the assumption and inductive hypothesis of this case, we know
{\footnotesize\begin{align*}
    & \forall \sigma \in \denot{\Gamma}.\forall \alpha_h \in \denot{\sigma(H)}.\forall \beta \in \denot{\Theta}. \exists \alpha \in \denot{\sigma(A)}.\exists \beta'. \msteptr{\alpha_h}{(\beta, e)}{\alpha}{(\beta', ())} \land \beta' \realz \sigma(F)
\end{align*}}\noindent
Since $\Gamma \vdash () : \nuot{unit}{\phi}$, we know $\sigma(\phi)$ holds.
Now, we need to show $\zassert{\phi}{e}$ can reduced into $e$ without add new effect, which is can be proved by \textsc{StAssert} and $\sigma(\phi)$. Then the proof immediate holds in this case.

\item:
{\footnotesize
\mprooftr{58mm}
{
$\Gamma \vdash \primop : t$ \quad $\Gamma \vdash t <: \overline{y{:}t}\sarr t_x$
$\forall i. \Gamma \vdash v_i : t_i$ \quad
$\Gamma, x{:}t_x\msubst{y}{v};\Delta;\Theta \vdash e : \rg{H}{A}{F}$
}
{TOpApp}
{$\Gamma;\Delta;\Theta \vdash \zlet{x{:}b}{\primop\ \overline{v}}{e} : \rg{H}{A}{F}$}
\ \\ \ \\}
This rule assume that $e \equiv \zlet{x{:}b}{\primop\ \overline{v}}{e}$, thus we need to prove
{\footnotesize\begin{align*}
\forall \sigma \in
    \denotation{\Gamma}. \forall \alpha_h \in \denot{\sigma(H)}.\forall \beta \in \denot{\Theta}. \exists \alpha \in \denot{\sigma(A)}.\exists \beta'. \msteptr{\alpha_h}{(\beta, \zlet{x{:}b}{\primop\ \overline{v}}{e})}{\alpha}{(\beta', ())} \land \beta' \realz \sigma(F)
\end{align*}}\noindent
From the assumption and inductive hypothesis of this case, we know
{\footnotesize\begin{alignat}{2}
\setcounter{equation}{0}
&\Delta(\eff{op}) = t \ptag{assumption} \label{real:3-1} \\
&\Gamma \vdash t <: \overline{y{:}t}\sarr t_x \ptag{assumption} \label{real:3-2} \\
&\forall i. \Gamma \vdash v_i : t_i \ptag{assumption} \label{real:3-3} \\
&\Gamma, x{:}t_x\msubst{y}{v};\Delta;\Theta \vdash e : \tau \ptag{assumption} \label{real:3-4} \\
&\forall v_x. \Gamma \vdash v_x : t_x\msubst{y}{v} \impl \Gamma;\Delta;\Theta \vdash e[x\mapsto v_x] : \tau[x\mapsto v_x] \ptag{Lemma~\ref{lemma:subt} and \ref{real:3-4}} \label{real:3-5} \\
& \forall \sigma \in \denot{\Gamma,  x{:}t_x\msubst{y}{v}}.\forall \alpha_h \in \denot{\sigma(H)}. \forall \beta \in \denot{\Theta}.  \notag
\\&\quad \exists \alpha \in \denot{\sigma(A)}. \exists \beta'. \msteptr{\alpha_h}{(\beta, e)}{\alpha}{(\beta', ())} \land \beta' \realz \sigma(F)
\ptag{induction hypothesis and \ref{real:3-4}} \label{real:3-6}
\end{alignat}}\noindent
This reduction step is pure, thus we can directly instantiate $\alpha$ in subgoal as $\alpha$ and apply hypothesis~\ref{real:3-6}, then which is sufficient to prove this case.

\item:
{\footnotesize
\mprooftr{38mm}
{$\Gamma;\Delta;\Theta \vdash e : \tau$ \quad
$\Gamma \vdash \tau <: \tau'$}
{TSub}
{$\Gamma;\Delta;\Theta \vdash e : \tau'$}
\ \\ \ \\}
The case can be directly proved by Lemma~\ref{lemma:subtyping}.

\end{enumerate}
\end{proof}

\paragraph{Fundamental Theorem} Now fundamental theorem can be proved with the help of Theorem~\ref{theorem:fundamental-ap} and Theorem~\ref{theorem:realizable}.

\begin{theorem}[Fundamental Theorem]\label{theorem:fundamental-original}A well-typed term, i.e., \(\Gamma; \Delta; \Theta \vdash e : \rg{H}{A}{F}\), generates traces consistent with the \PAT and can also terminate with the message buffer providing the capability.
{\small\begin{align*}
   &\forall \sigma \in
    \denotation{\Gamma}. \sigma(e) \in
    \denotation{\sigma(\rg{H}{A}{F})} \land  \forall \alpha_h \in \denot{\sigma(H)}.\forall \beta \in \denot{\Theta}. \exists \alpha.\exists \beta'. \msteptr{\alpha_h}{(\beta, e)}{\alpha}{(\beta', ())}
\end{align*}}
\end{theorem}
\begin{proof} For $\sigma \in\denotation{\Gamma}$, the first conjunct $\sigma(e) \in
    \denotation{\sigma(\rg{H}{A}{F})}$ can be provided directly via Theorem~\ref{theorem:fundamental-ap}. Additionally, for $ \alpha_h \in \denot{\sigma(H)}$ and  $\beta \in \denot{\Theta}$, Theorem~\ref{theorem:fundamental-ap} shows that $\exists \alpha \in \denot{\sigma(A)}.\exists \beta'. \msteptr{\alpha_h}{(\beta, e)}{\alpha}{(\beta', ())} \land \beta' \realz \sigma(F)$, which is sufficient to proved the second conjunct.
\end{proof}

\subsection{Type Soundness}

The type soundness can be proved by fundamental theorem and realizability.

\begin{theorem}[Type Soundness]\label{theorem:sound-ap} Given
  a well-formed handler specification $\Delta$, with ghost variables $\overline{x{:}b}$ and a violation property $A$, a controller $e$ that satisfies
  $\overline{x{:\rawnuot{b}{\top}}};\Delta; \emptyset \vdash e : \rg{\epsTL}{A}{\epsTL}$, then $e$ at least realize one trace consistent with $A$:
 {\footnotesize\begin{align*} \exists
    \overline{c{:}b}. \exists \alpha. \msteptr{[]}{(\emptyset,
      e[\overline{x\mapsto c}])}{\alpha}{(\emptyset, ())} \land \alpha \in
    \denot{A[\overline{x\mapsto c}]}
\end{align*}}
\end{theorem}
\begin{proof} According to the fundamental theorem, we have
{\footnotesize\begin{alignat}{2}
\setcounter{equation}{0}
    &\overline{x{:\rawnuot{b}{\top}}} ;\Delta; \emptyset \vdash e : \rg{\epsTL}{A}{\epsTL} \ptag{assumption} \label{sound:1}
    \\&\forall \sigma, \sigma \in
    \denot{\overline{x{:\rawnuot{b}{\top}}}} \impl \sigma(e) \in
    \denot{\sigma(\rg{\epsTL}{A}{\epsTL})} \ptag{\autoref{theorem:fundamental-ap} and \ref{sound:1}} \label{sound:2}
    \\&\forall \sigma, \sigma \in
    \denot{\overline{x{:\rawnuot{b}{\top}}}} \iff \exists \overline{c{:}b}. \sigma =\msubst{x}{c}  \ptag{definition of $\denot{\Gamma}$ and \ref{sound:2}}  \label{sound:3}
    \\&\forall \overline{c{:}b}. e\msubst{x}{c} \in
    \denot{\rg{\epsTL}{A\msubst{x}{c}}{\epsTL}} \ptag{hypothesis~\ref{sound:2} and \ref{sound:3}} \label{sound:4}
    \\&\forall \sigma \in
    \denotation{\Gamma}. \forall \alpha_h \in \denot{\sigma(\epsTL)}.\forall \beta \in \denot{\emptyset}. \notag
    \\&\qquad \exists \alpha \in \denot{\sigma(A)}.\exists \beta'. \msteptr{\alpha_h}{(\beta, e)}{\alpha}{(\beta', ())} \land \beta' \realz \sigma(\epsTL)
    \ptag{\autoref{theorem:realizable} and \ref{sound:1}} \label{sound:5}
    \\&\forall \alpha. \alpha \in \denot{\epsTL} \iff \alpha = []
    \ptag{Lemma~\ref{lemma:espTL}} \label{sound:6}
    \\&\forall \beta. \beta \in \denot{\emptyset} \iff \beta = \emptyset
    \ptag{Definition of capability denotation} \label{sound:7}
    \\&\forall \beta. \beta \realz \epsTL \iff \beta = \emptyset
    \ptag{Lemma~\ref{lemma:espTL-realz}} \label{sound:8}
    \\&\exists \overline{c{:}b}. \exists \alpha. \msteptr{[]}{(\emptyset, e\msubst{x}{c})}{\alpha}{(\emptyset, ())} 
    \ptag{\ref{sound:5} with \ref{sound:3},  \ref{sound:6}, \ref{sound:7}, \ref{sound:8}} \label{sound:9}
\end{alignat}}\noindent
Then, the $\alpha$ is the trace realized by the term $e$. Now we just need to prove $\alpha \in \denot{A\msubst{x}{c}}$.
Notice that the denotation of empty capability only contains empty buffer, 
the definition of \PAT denotation as shown in \autoref{fig:type-denotation} indicates
{\footnotesize\begin{align*}
    &\forall \alpha_h \in\denot{\sigma(\epsTL)}.\forall \alpha_f \in\denot{\sigma(\epsTL)}. \forall \alpha\; \beta\;\beta'\;e_h\;e_f.
  \\&\qquad \msteptr{[]}{(\emptyset, e_h)}{\alpha_h}{(\beta, ())}\land \msteptr{\alpha_h}{(\beta, e)}{\alpha}{(\beta', ())} \land \msteptr{\alpha_h \listconcat \alpha }{(\beta', e_f)}{\alpha_f}{(\emptyset, ())} \impl \alpha \in\denot{A\msubst{x}{c}}
\end{align*}}\noindent
Again, according to Lemma~\ref{lemma:espTL},
{\footnotesize\begin{alignat}{2}
  &\msteptr{[]}{(\emptyset, ())}{[]}{(\emptyset, ())} \ptag{definition of $\hookrightarrow^*$} \label{sound:10}
  \\&\msteptr{\alpha}{(\emptyset, ())}{[]}{(\emptyset, ())} \ptag{definition of $\hookrightarrow^*$} \label{sound:11}
  \\&\exists \overline{c{:}b}.\exists \alpha \in\denot{\sigma(A\msubst{x}{c})}. \msteptr{[]}{(\emptyset, e)}{\alpha}{(\emptyset, ())}
\ptag{Denotation of \PAT, \ref{sound:9},\ref{sound:10}, and \ref{sound:11}}
  \label{sound:12}
\end{alignat}}\noindent
This is sufficient to establish the original theorem we aim to prove.
\end{proof}

% Our synthesis algorithm is sound, ensuring that it always produces type-safe controllers. Furthermore, the synthesized controller can realize an execution trace when the handler specification $\Delta$ is complete, meaning that all behaviors permitted by $\Delta$ are realizable.

% \begin{theorem}\label{theorem:algo-well-typed}[Well-type-ness]
% The controller synthesized by the algorithm can also be derived using our declarative rules.
% \end{theorem}

% \begin{definition}[Complete Handler Specification] A well-formed handler specification $\Delta$ is complete iff for all $\Delta(\eff{op}) = \overline{x{:}t} \sarr \rg{H}{A}{\bigwedge_{i = 1 ... n} \finalA \msg{op_i}{\phi_i}}$,
% {\small\begin{align*}
%     \forall \alpha \in \denot{A}. \forall \overline{c \in \denot{t}}. \forall \beta. \beta \in \denot{\seqA \overline{\lastA \msg{op_i}{\phi_i}}} \impl \alpha \vDash \zevent{op}{\overline{c}} \Downarrow \beta
% \end{align*}}
% \end{definition}



% \begin{theorem}\label{theorem:algo-sound}[Algorithm Soundness] Under a complete handler specification, with ghost variables $\overline{x{:}b}$ and a safety property $A$, the synthesized program $e$ realizes at least one trace that violate $A$, i.e.,
% {\small\begin{align*}
% \exists \overline{c{:}b}. \exists \alpha.
% \msteptr{[]}{([],e[\overline{x\mapsto c}])}{\alpha}{([], ())} \land \alpha \not\in
%     \denot{A[\overline{x\mapsto c}]}
% \end{align*}}
% \end{theorem}
